\documentclass[10pt,a4paper]{amsart}
\usepackage[margin=19mm]{geometry}
\usepackage{amsmath,amssymb,mathtools}
\usepackage[scr=rsfs,cal=euler]{mathalfa}
\usepackage{xfrac}
\usepackage[unicode,hidelinks,bookmarks=false]{hyperref}
\newcommand{\ie}{i.e.\ }
\newcommand{\cf}{cf.\ }
\newcommand{\resp}{respectively\ }
\newcommand{\Subsection}{\subsection}
\providecommand{\nobreakdash}{\nobreak\hbox{-}}
\setlength{\emergencystretch}{2em}
\multlinegap0pt
\usepackage{amssymb}
\usepackage{mathtools}
\usepackage{braket}
\newtheorem{prop}{Proposition}
\DeclareMathOperator{\Par}{Par}
\newcommand{\bN}{\mathbb{N}}
\newcommand{\bfP}{\mathbf{P}}
\newcommand{\bfb}{\mathbf{b}}
\newcommand{\cE}{\mathcal{E}}
\DeclareMathOperator{\length}{length}
\DeclareMathOperator{\oF}{F}
\newcommand{\pf}{\mathrm{pf}}
\newcommand{\sG}{\mathscr{G}}
\newcommand{\sP}{\mathscr{P}}
\title{The supersingular locus of the Shimura variety of $\mathrm{GU}(2,n-2)$}
\author{Maria Fox, Naoki Imai}
\date{Source: arXiv:2108.03584v3 (CC BY 4.0)}
\begin{document}
\maketitle

\noindent\emph{Source and licence.} The supersingular locus of the Shimura variety of $\mathrm{GU}(2,n-2)$. Version arXiv:2108.03584v3 (CC BY 4.0), distributed by arXiv under the Creative Commons Attribution 4.0 International licence, http://creativecommons.org/licenses/by/4.0/, as recorded in the arXiv licence metadata for this exact version and shown on the abstract page https://arxiv.org/abs/2108.03584v3. Full licence notice: \texttt{licenses/arxiv-2108.03584-CC-BY-4.0.txt}.\par\smallskip

\noindent\textit{Editorial scope.} Counted unit: the article's prop prop:inter (source0.tex lines 2240--2280). The article's complete original proof of it is delivered with it (source0.tex lines 2281--2391, one \texttt{proof} environment of 111 lines). No mathematics is written, repaired or re-derived: the statement and the proof are the article's own lines, in the article's own notation, and no retained line is substituted or re-flowed.  No further source-local context is needed: every cross-reference of every retained line resolves inside the two delivered spans.  Nothing was omitted from any retained span, so the package carries no omission record.  No retained line contains a citation, so the wrapper carries no bibliography and no bibliography entry.  No retained line contains a figure environment, an image-inclusion command or drawing code, so no image is delivered and the package declares no asset dependency.  Editorial formatting only, in addition: an AMS \texttt{amsart} preamble that restates the article's own declarations and macros used by the retained lines, namely the environments \texttt{prop} on separate counters, together with the standard packages those lines need.  The retained environments print in this document as Proposition 1 in order of appearance, whereas the article prints the same items with the numbering of its own full document.  The complete native source of the article as distributed in the arXiv e-print for this exact version is bundled unchanged as \texttt{sources/114750/source0.tex}.
\begin{prop}\label{prop:inter}
Assume that 
$\mathring{X}_{\mu^*}^{\bfb_i,x}(\tau_i^*) \cap \mathring{X}_{\mu^*}^{\bfb_{i'},x'}(\tau_{i'}^*)$ 
is not empty. 
Then we have $l_{x,x'} \leq i+i'-1$. 

The subscheme 
$\mathring{X}_{\mu^*}^{\bfb_i,x}(\tau_i^*) \cap \mathring{X}_{\mu^*}^{\bfb_{i'},x'}(\tau_{i'}^*) \subset \mathring{X}_{\mu^*}^{\bfb_{i},x}(\tau_{i}^*)$ is the locus defined by the condition that 
$\cE_x +\varpi \cE_{x'} \subset \cE_+ \subset \frac{1}{\varpi} \cE_{x,x'}^-$, 
$\varpi \cE_{x,x'}^+ \subset \cE_- \subset \cE_x \cap \frac{1}{\varpi} \cE_{x'}$, 
\begin{equation}\label{eq:lengE-}
 \length ((\cE_- +\cE_{x,x'}^-)/ \cE_{x,x'}^- ) \leq  [(i'-i+d_2 -d_1 )/2] 
\end{equation} 
and 
\begin{equation}\label{eq:lengsum}
 \length ((\cE +\cE_{+,-} )/ \cE_{+,-}) + \length ((\cE_- +\cE_{x,x'}^-)/ \cE_{x,x'}^- ) = i'-1 -d_1. 
\end{equation} 
In particular,  
$\mathring{X}_{\mu^*}^{\bfb_i,x}(\tau_i^*) \cap \mathring{X}_{\mu^*}^{\bfb_{i'},x'}(\tau_{i'}^*)$ 
is the union of 
$\left( \mathring{X}_{\mu^*}^{\bfb_i,x}(\tau_i^*) \cap \mathring{X}_{\mu^*}^{\bfb_{i'},x'}(\tau_{i'}^*) \right)_{\bfP_{x,x'},[w_{j_1,j_2}]}$ 
for 
$j_1,j_2 \in \bN$ such that $j_1 +d_2 -i \leq j_2 \leq j_1 +d_2-i +1$, 
\begin{align*}
    &i-1-d_2 \leq j_1 \leq i-1-d_1
    ,\\  
 &
 i' -i -d_1 \leq 
 j_2 \leq \min \{ 
 [(i'-i+d_2 -d_1 )/2] 
 , n-i-d_2 -j_1 
 \}. 
\end{align*}
Further we have 
\[
 \left( \mathring{X}_{\mu^*}^{\bfb_i,x}(\tau_i^*) \cap \mathring{X}_{\mu^*}^{\bfb_{i'},x'}(\tau_{i'}^*)  \right)_{\bfP_{x,x'},[w_{j_1,j_2}]} = 
 \mathring{X}_{\mu^*}^{\bfb_i,x}(\tau_i^*)_{\bfP_{x,x'},[w_{j_1,j_2}]} 
 \cap \Par_{t_{j_1,j_2}} 
(\sG_{j_1,j_2}; \sP_{j_1,j_2})_{r_{j_1,j_2}}^{\pf} . 
\] 
\end{prop}

\begin{proof}
The intersection 
$\mathring{X}_{\mu^*}^{\bfb_i,x}(\tau_i^*) \cap
 \mathring{X}_{\mu^*}^{\bfb_{i'},x'}(\tau_{i'}^*)$ 
is parametrized by 
$\cE \dashrightarrow \cE_x$ which is equal to 
$\nu_i^*$ such that 
$\cE \subset \oF (\cE^{\vee}) \subset \frac{1}{\varpi} \cE$ 
and 
$\cE \dashrightarrow \cE_{x'}$ 
is equal to $\nu_{i'}^*$. 
Let $\cE$ be a point of 
$\mathring{X}_{\mu^*}^{\bfb_i,x}(\tau_i^*) \cap \mathring{X}_{\mu^*}^{\bfb_{i'},x'}(\tau_{i'}^*)$. 
We put 
\[
 \cE_+=\cE +\cE_x, \quad 
 \cE_-=\cE \cap \cE_x, \quad 
 \cE_+'=\cE +\cE_{x'}, \quad
 \cE_-'=\cE \cap \cE_{x'}. 
\]
Then we have 
\begin{align*}
&\length (\cE_x/\cE_-)=i, \quad 
\length (\cE/\cE_-)=i-1, \\ 
&\length (\cE_{x'}/\cE_-')= i', \quad 
\length (\cE/\cE_-') = i' -1. 
\end{align*}
Hence we have 
$\length(\cE_-/(\cE_- \cap \cE_-')) \leq i' -1$ and 
$\length (\cE_x/(\cE_- \cap \cE_-')) \leq i+i' -1$. 
Therefore the inclusion 
$\cE_- \cap \cE_-' \subset \cE_x \cap \cE_{x'}$ implies that 
\[
 l_{x,x'} \leq i+i' -1. 
\]

We have 
$\cE_x +\varpi \cE_{x'} \subset \cE_+$ and 
$\varpi \cE_{x,x'}^+ =\varpi \cE_x +(\cE_x \cap \varpi \cE_{x'}) \subset \cE_-$, 
since $\varpi \cE_{x'} \subset \cE$. 
We have 
$\varpi \cE_+ \subset \cE_x \cap (\cE_{x'}+\varpi \cE_x)=\cE_{x,x'}^-$ and  
$\cE_- \subset \cE_x \cap \frac{1}{\varpi} \cE_{x'}$, 
since $\varpi \cE_+ \subset \cE_x$ and 
$\varpi \cE \subset \cE_{x'}$. 

We put $j_1=\length ((\cE_+ +\cE_{x,x'}^+ )/ \cE_{x,x'}^+ )$ 
and $j_2=\length ((\cE_- +\cE_{x,x'}^- )/\cE_{x,x'}^- )$. 
We have 
\begin{align*}
 \length (\cE_- / (\cE_- \cap \cE_-')) 
 &=\length (\cE_- / (\cE_- \cap \cE_{x'})) 
 =j_2 +\length ((\cE_-\cap \cE_{x,x'}^-) / (\cE_- \cap \cE_{x'}))\\ 
 &=j_2 +\length (\cE_{x,x'}^-/ (\cE_x \cap \cE_{x'})) = j_2 +d_1. 
\end{align*}
We have 
\[
 j_2 +i-d_2 = \length ((\cE_- + \cE_{x,x'}^-)/\cE_-) 
 \leq 1+\length ((\oF (\cE_+^{\vee}) + \cE_{x,x'}^-)/\oF (\cE_+^{\vee}))
\] 
since $\length (\oF (\cE_+^{\vee})/\cE_-)=1$. 
Further we have 
\begin{align*}
 \length & ((\oF (\cE_+^{\vee}) + \cE_{x,x'}^-)/\oF (\cE_+^{\vee}))
 = \length (\cE_+ / (\cE_+ \cap \cE_{x,x'}^+)) 
 \leq \length (\cE_+ / (\cE_-' +\cE_x )) \\ 
 &\leq 
 \length (\cE / (\cE_- +\cE_-'))
 =  i'-1-\length (\cE_- / (\cE_- \cap \cE_-')) 
 =i'-1-j_2 -d_1. 
\end{align*}
Therefore we obtain 
$j_2  \leq  [(i'-i+d_2 -d_1 )/2]$. 

Further, 
$j_1 +d_2 -i \leq j_2 \leq j_1 +d_2-i +1$ 
follows from $\length (\oF (\cE_-^{\vee})/\cE_+)=1$. 
This implies 
$i-1-d_2 \leq j_1$ and 
$d_2 -i \leq j_2$. 
We have 
$j_1 \leq i-1-d_1$ and 
$j_2 \leq n-i-d_2 -j_1$ 
by the inclusions 
$\cE_x +\varpi \cE_{x'} \subset \cE_+ \cap \cE_{x,x'}^+$
and 
$\cE_+ +\cE_{x,x'}^+ \subset \cE_- \cap \cE_{x,x'}^-$. 
The equality 
\[
 \length ((\cE +\cE_{+,-} )/ \cE_{+,-}) + \length ((\cE_- +\cE_{x,x'}^-)/ \cE_{x,x'}^- ) = i'-1 -d_1 
\] 
and $\length ((\cE +\cE_{+,-} )/ \cE_{+,-}) \leq i-1$ imply 
$j_2 \geq i'-i-d_1$. 

We have 
\begin{align*}
 \length & ((\cE +\cE_{+,-} )/ \cE_{+,-}) 
 = \length (\cE / (\cE \cap \cE_{+,-}))
 = \length (\cE / (\cE \cap (\cE_- + \cE_{x'} ))) \\ 
 &= \length ((\cE +\cE_{x'})/ (\cE_- + \cE_{x'} )) 
 = \length ((\cE +\cE_{x'})/ \cE_{x'} ) 
 -\length ((\cE_- + \cE_{x'} )/ \cE_{x'} ) \\ 
 &= \length ((\cE +\cE_{x'})/ \cE_{x'} ) 
 -\length (\cE_-/ (\cE_- \cap \cE_{x'} ))  . 
\end{align*}
Hence, 
$\length  ((\cE +\cE_{+,-} )/ \cE_{+,-}) = i'-1 - j_2 -d_1$ 
if and only if 
$\length ((\cE +\cE_{x'})/ \cE_{x'} ) = i'-1$. 
This implies the last claim. 
\end{proof}

\end{document}
